\begin{itemize}
	\item \textbf{生理因素}：激素水平变化（围绝经期、产后、哺乳期）、慢性疾病（甲状腺功能异常、糖尿病、贫血）、药物影响（抗抑郁药、部分降压药、激素类避孕药）、疲劳与睡眠不足。
	\item \textbf{心理因素}：压力、焦虑、抑郁、性内疚或羞耻、既往性创伤、对身体的负面评价、注意力长期被工作或育儿占据。
	\item \textbf{关系因素}：情感疏离、沟通不畅、未解决的冲突、伴侣的性技巧或态度令其不适、长期"性生活的脚本"单调乏味。
	\item \textbf{情境因素}：缺乏私密空间与时间、孩子同床、身心疲惫、把性当成"任务"而非享受。
	\item \textbf{分化与评估}：需区分\textbf{对所有人、所有情境都缺乏性欲}（广义性欲低下）与\textbf{只对伴侣或特定情境缺乏性欲}（情境性/关系性）；后者提示关系与情境因素更关键，前者需更多排查生理因素。
	\item \textbf{应对}：先处理可逆因素（改善睡眠、减压、治疗基础疾病、与医生复核用药）；其次改善关系与沟通，把"性"从任务重新变回连接；必要时寻求性治疗师或妇科/内分泌科评估，明确是否存在 HSDD 等情况并考虑相应治疗（详见下文 HSDD 一节）。
\end{itemize}

\noindent\textit{【编者注】}性欲低下是女性最常见的性困扰之一，\textbf{多数并非"冷淡"，而是可以被理解和改善的}。若长期为此困扰，或已影响关系与情绪，请寻求专业帮助——它与"爱不爱对方"无关。
